% Extracted from 群论练习题-20240605(1).pdf.

\providecommand{\rank}{\operatorname{rank}}
\providecommand{\Sym}{\operatorname{Sym}}
\providecommand{\Syl}{\operatorname{Syl}}
\providecommand{\Hall}{\operatorname{Hall}}
\providecommand{\Hol}{\operatorname{Hol}}
\providecommand{\Inn}{\operatorname{Inn}}
\providecommand{\Out}{\operatorname{Out}}
\providecommand{\Fit}{\operatorname{Fit}}
\providecommand{\Core}{\operatorname{Core}}

\begin{exercise}
Find conjugacy classes and fusion classes of \(A_5\).
\end{exercise}

\begin{proof}
The conjugacy classes of \(A_5\) are obtained by intersecting the \(S_5\)
cycle-type classes with \(A_5\). The result is
\[
    \{1\},
\]
the \(15\) double transpositions,
\[
    (ab)(cd),
\]
the \(20\) \(3\)-cycles, and two classes of \(5\)-cycles, each of size \(12\).
For example, if \(x=(1\,2\,3\,4\,5)\), then the two \(A_5\)-classes of
\(5\)-cycles are represented by \(x\) and \(x^2\).

The fusion classes are the \(\Aut(A_5)\)-orbits. Since
\[
    \Aut(A_5)\cong S_5,
\]
these are just the \(S_5\)-conjugacy classes contained in \(A_5\):
\[
    \{1\},\qquad
    \{x\in A_5: |x|=2\},\qquad
    \{x\in A_5: |x|=3\},\qquad
    \{x\in A_5: |x|=5\}.
\]
Thus the two \(A_5\)-classes of \(5\)-cycles fuse under \(\Aut(A_5)\).
\end{proof}

\begin{exercise}
    Determine the conjugacy classes and fusion classes of
    \[
        (\mathbb Z/p^e\mathbb Z)^d
        \qquad\text{and}\qquad
        \mathbb Z/p^a\mathbb Z\times \mathbb Z/p^b\mathbb Z.
    \]
    \end{exercise}
    
    \begin{proof}
    Both groups are abelian. Hence every conjugacy class is a singleton.
    Therefore it remains only to determine the fusion classes, that is, the
    orbits under the automorphism group.
    
    First let
    \[
        G=(\mathbb Z/p^e\mathbb Z)^d.
    \]
    Put
    \[
        R=\mathbb Z/p^e\mathbb Z.
    \]
    Then \(G\) is a free \(R\)-module of rank \(d\), and
    \[
        \Aut(G)=\GL_d(R).
    \]
    
    For \(0\neq v\in G\), define
    \[
        \nu(v)=\max\{i\mid v\in p^iG\}.
    \]
    Equivalently,
    \[
        v\in p^iG\setminus p^{i+1}G
        \quad\Longleftrightarrow\quad
        \nu(v)=i.
    \]
    Since each subgroup \(p^iG\) is characteristic in \(G\), the integer
    \(\nu(v)\) is invariant under \(\Aut(G)\).
    
    Conversely, suppose
    \[
        v\in p^iG\setminus p^{i+1}G.
    \]
    Then
    \[
        v=p^i u
    \]
    for some \(u\in G\setminus pG\). Since \(R\) is local with maximal ideal
    \((p)\), the condition \(u\notin pG\) means that \(u\) is a primitive vector.
    Hence \(u\) can be extended to an \(R\)-basis of \(G\). Therefore some element
    of \(\GL_d(R)\) sends \(u\) to \(e_1\), and hence sends \(v\) to
    \[
        p^i e_1.
    \]
    Thus the fusion classes of \(G=(\mathbb Z/p^e\mathbb Z)^d\) are
    \[
        \{0\},
        \qquad
        p^iG\setminus p^{i+1}G
        \quad (0\leq i<e).
    \]
    Since an element of \(p^iG\setminus p^{i+1}G\) has order \(p^{e-i}\), this is
    equivalently the statement that two nonzero elements are fused if and only if
    they have the same order.
    
    Now let
    \[
        G=\langle x\rangle\oplus \langle y\rangle
        \cong C_{p^a}\times C_{p^b},
        \qquad a\geq b,
    \]
    where
    \[
        |x|=p^a,
        \qquad
        |y|=p^b.
    \]
    We shall write
    \[
        p^a x=0,
        \qquad
        p^b y=0.
    \]
    
    We claim that the fusion classes are represented uniquely by the elements
    \[
        p^i x+p^j y,
    \]
    where
    \[
        0\leq i\leq a,
        \qquad
        0\leq j\leq b,
    \]
    and the pair \((i,j)\) satisfies the canonical conditions
    \[
        j\leq i,
        \qquad
        b-j\leq a-i.
    \]
    Equivalently,
    \[
        j\leq i\leq a-b+j.
    \]
    Thus the fusion classes are precisely
    \[
        \Aut(G)\cdot (p^i x+p^j y)
    \]
    for all such canonical pairs \((i,j)\).
    
    We prove existence first. We shall use the following elementary automorphisms:
    \[
        x\mapsto ux,\qquad y\mapsto vy,
    \]
    where \(u\in(\mathbb Z/p^a\mathbb Z)^\times\) and
    \(v\in(\mathbb Z/p^b\mathbb Z)^\times\), and also
    \[
        x\mapsto x+cy,\qquad y\mapsto y,
    \]
    as well as
    \[
        x\mapsto x,\qquad y\mapsto y+p^{a-b}dx.
    \]
    These are automorphisms because their inverses are obtained by replacing
    \(c\) or \(d\) by \(-c\) or \(-d\), and because \(p^{a-b}x\) has order
    dividing \(p^b\).
    
    Let
    \[
        g=\alpha x+\beta y\in G.
    \]
    If \(g=0\), then
    \[
        g=p^a x+p^b y,
    \]
    which corresponds to the canonical pair \((a,b)\).
    
    Assume now that \(g\neq 0\). Define
    \[
        r=v_p(\alpha),
        \qquad
        s=v_p(\beta),
    \]
    with the conventions
    \[
        r=a \quad\text{if } \alpha=0,
        \qquad
        s=b \quad\text{if } \beta=0.
    \]
    Then the order of \(g\) is
    \[
        |g|=p^m,
        \qquad
        m=\max\{a-r,\ b-s\}.
    \]
    
    There are two cases.
    
    First suppose
    \[
        a-r\geq b-s.
    \]
    Then the \(x\)-component has maximal order. In this case \(r<a\). By scaling
    \(x\), we may assume that the \(x\)-coefficient is exactly \(p^r\). Thus
    \[
        g\sim p^r x+\beta y.
    \]
    If \(s<r\), then either \(\beta=0\), in which case \(s=b\), or we may scale
    \(y\) so that \(\beta=p^s\). Hence
    \[
        g\sim p^r x+p^s y.
    \]
    The pair \((r,s)\) is canonical, since
    \[
        s<r
    \]
    and
    \[
        b-s\leq a-r
    \]
    by assumption.
    
    If \(s\geq r\), then \(r\leq b\). If \(r=b\), then \(\beta=0\), and
    \[
        g\sim p^r x=p^r x+p^b y.
    \]
    The pair \((r,b)\) is canonical. If \(r<b\), write
    \[
        \beta=p^r w
        \qquad
        \text{in } \mathbb Z/p^b\mathbb Z.
    \]
    Using the automorphism
    \[
        x\mapsto x+cy,
        \qquad
        y\mapsto y,
    \]
    we may replace the \(y\)-coefficient by
    \[
        p^r(w+c).
    \]
    Choosing \(c\) so that
    \[
        w+c\equiv 1\pmod{p^{b-r}},
    \]
    we obtain
    \[
        g\sim p^r x+p^r y.
    \]
    The pair \((r,r)\) is canonical.
    
    Now suppose
    \[
        b-s>a-r.
    \]
    Then the \(y\)-component has maximal order, so \(s<b\). Put
    \[
        i=a-b+s.
    \]
    The inequality
    \[
        b-s>a-r
    \]
    is equivalent to
    \[
        r>i.
    \]
    After scaling \(y\), we may assume
    \[
        g=\alpha x+p^s y.
    \]
    Using the automorphism
    \[
        x\mapsto x,
        \qquad
        y\mapsto y+p^{a-b}x,
    \]
    we obtain
    \[
        g\mapsto
        \bigl(\alpha+p^{s+a-b}\bigr)x+p^s y
        =
        \bigl(\alpha+p^i\bigr)x+p^s y.
    \]
    Since \(r>i\), the element \(\alpha\) is divisible by \(p^{i+1}\). Hence
    \[
        \alpha+p^i=p^i u
    \]
    for some unit \(u\). Scaling \(x\), we get
    \[
        g\sim p^i x+p^s y.
    \]
    The pair \((i,s)\) is canonical because
    \[
        s\leq a-b+s=i
    \]
    and
    \[
        b-s=a-i.
    \]
    This proves that every element is automorphic to some canonical representative.
    
    It remains to prove uniqueness. For \(g\in G\), define its height by
    \[
        h(g)=\max\{t\mid g\in p^tG\},
    \]
    with the convention \(h(0)=\infty\). Since each \(p^tG\) is characteristic,
    height is invariant under automorphisms. Of course, the order of an element is
    also invariant under automorphisms.
    
    Let
    \[
        g_{i,j}=p^i x+p^j y
    \]
    be a canonical representative. The canonical condition
    \[
        b-j\leq a-i
    \]
    implies
    \[
        |g_{i,j}|=p^{a-i}.
    \]
    Therefore the order of \(g_{i,j}\) determines \(i\).
    
    For fixed \(i\), the height of \(g_{i,j}\) is
    \[
        h(g_{i,j})=
        \begin{cases}
            j, & j<b,\\
            i, & j=b.
        \end{cases}
    \]
    This determines \(j\) among the canonical possibilities. Indeed, if \(j<b\), then
    \(j=h(g_{i,j})\). If \(j=b\), then the condition \(j\leq i\) gives \(b\leq i\).
    Thus no canonical representative with \(j'<b\) can have the same height \(i\),
    because \(j'=i\) would force \(i<b\), contradicting \(b\leq i\).
    
    Hence two canonical representatives are automorphic only if their pairs
    \((i,j)\) are equal. This proves uniqueness.
    
    Therefore, for
    \[
        G\cong C_{p^a}\times C_{p^b},
        \qquad a\geq b,
    \]
    the fusion classes are represented uniquely by
    \[
        p^i x+p^j y
    \]
    with
    \[
        0\leq i\leq a,
        \qquad
        0\leq j\leq b,
        \qquad
        j\leq i,
        \qquad
        b-j\leq a-i.
    \]
    If \(b>a\), the same description applies after interchanging the two cyclic
    factors.
    \end{proof}

\begin{exercise}
Let \(G = \mathrm{D}_{2n}\), a dihedral group of order \(2n\). Find
\(\operatorname{Aut}(G)\), and determine the conjugacy classes and fusion classes of
elements of \(G\).
\end{exercise}

\begin{proof}
Write
\[
    G=\langle r,s\mid r^n=s^2=1,\ srs=r^{-1}\rangle.
\]
The cyclic subgroup \(\langle r\rangle\) is characteristic, since it is the
unique cyclic subgroup of index \(2\) except in the small cases where the same
formula is checked directly. Hence an automorphism has the form
\[
    r\mapsto r^u,\qquad s\mapsto r^v s,
\]
where \(u\in(\mathbb Z/n\mathbb Z)^\times\) and \(v\in\mathbb Z/n\mathbb Z\).
Conversely, every such choice defines an automorphism. Therefore
\[
    \Aut(D_{2n})\cong \mathbb Z_n\rtimes (\mathbb Z/n\mathbb Z)^\times,
\]
the affine group on \(\mathbb Z_n\).

The conjugacy classes of rotations are
\[
    \{r^i,r^{-i}\},
\]
with the obvious singleton exceptions \(r^0=1\), and \(r^{n/2}\) when \(n\) is
even. If \(n\) is odd, all reflections lie in one conjugacy class:
\[
    \{r^i s\mid 0\leq i<n\}.
\]
If \(n\) is even, the reflections split into two conjugacy classes:
\[
    \{r^{2i}s\mid 0\leq i<n/2\},
    \qquad
    \{r^{2i+1}s\mid 0\leq i<n/2\}.
\]

For fusion classes, the automorphism group sends \(r^i\) to \(r^{ui}\), with
\(u\) a unit modulo \(n\). Thus the rotation fusion classes are determined by
\(\gcd(i,n)\), equivalently by the order of \(r^i\). The formula
\[
    s\mapsto r^v s
\]
shows that all reflections are fused by \(\Aut(G)\). Hence the fusion classes
are the identity, the sets of rotations of each fixed order \(>1\), and one
class consisting of all reflections.
\end{proof}

\begin{exercise}
Let \(G\) be a \(p\)-group, where \(p\) be a prime.
\begin{enumerate}[label=(\alph*)]
    \item Prove that the center \(Z(G) \neq 1\).
    \item Classify groups of order \(p^3\).
\end{enumerate}
\end{exercise}

\begin{proof}
\((a)\) The class equation gives
\[
    |G|
    =
    |Z(G)|+\sum_i [G:C_G(x_i)],
\]
where \(x_i\) runs over representatives of the noncentral conjugacy classes.
Each noncentral class has size divisible by \(p\). Hence
\[
    |Z(G)|\equiv |G|\equiv 0\pmod p.
\]
Thus \(p\mid |Z(G)|\), and so \(Z(G)\neq 1\).

\((b)\) If \(|G|=p^3\) and \(G\) is abelian, then
\[
    G\cong C_{p^3},\qquad
    C_{p^2}\times C_p,\qquad
    C_p\times C_p\times C_p.
\]
Now suppose that \(G\) is nonabelian. Since \(G\) is a \(p\)-group,
\(Z(G)\neq 1\). If \(|Z(G)|=p^2\), then \(G/Z(G)\) is cyclic, forcing \(G\) to
be abelian. Therefore
\[
    |Z(G)|=p,
    \qquad
    G'=Z(G)\cong C_p.
\]
For \(p=2\), the two nonabelian possibilities are
\[
    D_8,\qquad Q_8.
\]
For odd \(p\), the two nonabelian possibilities are
\[
    \langle a,b\mid a^{p^2}=b^p=1,\ a^b=a^{1+p}\rangle
\]
and
\[
    \langle a,b,c
        \mid
        a^p=b^p=c^p=1,\ c=[a,b],\ [a,c]=[b,c]=1
    \rangle.
\]
The first has exponent \(p^2\), and the second has exponent \(p\).
\end{proof}

\begin{exercise}
Let
\[
    G = \langle a,b \mid |a| = 2^e,\ |b| = 2,\ a^b = a^{2^{e-1}-1} \rangle,
\]
where \(e \geq 4\). Prove that \(G\) contains a subgroup which is isomorphic to
\(\mathbb{Z}_{2^e}\), \(\mathrm{D}_{2^e}\) and \(\mathrm{Q}_{2^e}\).
\end{exercise}

\begin{proof}
The subgroup
\[
    \langle a\rangle
\]
is cyclic of order \(2^e\), so it is isomorphic to \(\mathbb Z_{2^e}\).

Put \(c=a^2\). Then \(|c|=2^{e-1}\), and
\[
    c^b=(a^b)^2=a^{2(2^{e-1}-1)}=a^{-2}=c^{-1}.
\]
Hence
\[
    \langle c,b\rangle
    =
    \langle c,b\mid c^{2^{e-1}}=b^2=1,\ c^b=c^{-1}\rangle
    \cong D_{2^e}.
\]

Finally put \(d=ab\). Then
\[
    d^2=(ab)^2=aa^b=a^{2^{e-1}},
\]
and
\[
    c^d=c^b=c^{-1}.
\]
Since
\[
    d^2=a^{2^{e-1}}=c^{2^{e-2}},
\]
we get
\[
    \langle c,d\rangle
    =
    \langle c,d\mid c^{2^{e-1}}=1,\ d^2=c^{2^{e-2}},\ c^d=c^{-1}\rangle
    \cong Q_{2^e}.
\]
\end{proof}

\begin{exercise}
Let \(G\) be a finite group of even order. Show that \(G\) has an odd number of
elements of order \(2\).
\end{exercise}

\begin{proof}
Pair each element \(x\in G\) with its inverse \(x^{-1}\). The elements not fixed
by this pairing occur in pairs. The fixed elements are exactly those satisfying
\[
    x=x^{-1},
\]
namely \(x^2=1\). Since \(|G|\) is even, the number of solutions of \(x^2=1\)
is even. One of them is \(1\), so the number of nonidentity solutions of
\(x^2=1\), that is the number of elements of order \(2\), is odd.
\end{proof}

\begin{exercise}
Prove that \(A_n\) with \(n \geq 5\) is a simple group.
\end{exercise}

\begin{proof}
Let \(1\neq N\trianglelefteq A_n\). Choose \(1\neq x\in N\) with the number of
moved points minimal. If \(x\) moves more than \(3\) points, then one can choose
a \(3\)-cycle \(t\in A_n\) such that the commutator
\[
    [x,t]=x^{-1}t^{-1}xt
\]
is nontrivial and moves fewer points than \(x\), contradicting the choice of
\(x\). Hence \(N\) contains an element moving exactly \(3\) points, so \(N\)
contains a \(3\)-cycle.

All \(3\)-cycles are conjugate in \(A_n\) for \(n\geq 5\). Since \(N\) is normal,
\(N\) contains every \(3\)-cycle. But \(A_n\) is generated by its \(3\)-cycles.
Therefore
\[
    N=A_n.
\]
Thus \(A_n\) is simple.
\end{proof}

\begin{exercise}
Let \(G\) be a cyclic group of order \(n\). Determine \(\operatorname{Aut}(G)\).
\end{exercise}

\begin{proof}
Let \(G=\langle g\rangle\). An automorphism is determined by the image of \(g\),
and \(g^i\) is again a generator if and only if
\[
    \gcd(i,n)=1.
\]
Thus
\[
    \Aut(G)\cong (\mathbb Z/n\mathbb Z)^\times.
\]
Under this isomorphism the unit \(\overline{i}\) corresponds to the automorphism
\[
    g\mapsto g^i.
\]
\end{proof}

\begin{exercise}
Let \(G = \mathbb{Z}_p^d\), where \(p\) is a prime and \(d\) is a positive
integer. Then \(\operatorname{Aut}(G) \cong \mathrm{GL}(d,p)\).
\end{exercise}

\begin{proof}
The group \(G\) is the additive group of the vector space
\[
    V=\mathbb F_p^d.
\]
Every group homomorphism \(G\to G\) preserves addition, hence is an
\(\mathbb F_p\)-linear map. It is an automorphism precisely when the
corresponding linear map is invertible. Therefore
\[
    \Aut(G)=\GL(V)\cong \mathrm{GL}(d,p).
\]
\end{proof}

\begin{exercise}
Let \(G = T^d\), where \(T\) is a nonabelian simple group. Then
\[
    \operatorname{Aut}(G) = \operatorname{Aut}(T) \wr S_d,
\]
and
\[
    \operatorname{Hol}(G)
    =
    \operatorname{Hol}(T^d)
    =
    (T \rtimes \operatorname{Aut}(T)) \wr S_d
    =
    (\check{G} \times \hat{G}).(\operatorname{Out}(T) \wr S_d).
\]
In particular, if \(G = T\), then
\[
    \operatorname{Hol}(T)
    =
    T \rtimes \operatorname{Aut}(T)
    =
    (\check{T} \times \hat{T}).\operatorname{Out}(T).
\]
\end{exercise}

\begin{proof}
Write
\[
    G=T_1\times \cdots\times T_d,
    \qquad T_i\cong T.
\]
The subgroups \(T_i\) are exactly the minimal normal subgroups of \(G\). Hence
every automorphism of \(G\) permutes the set
\[
    \{T_1,\ldots,T_d\}.
\]
The kernel of this permutation action is
\[
    \Aut(T_1)\times\cdots\times \Aut(T_d)\cong \Aut(T)^d,
\]
and every permutation of the \(d\) factors is induced by an automorphism of
\(G\). Thus
\[
    \Aut(G)\cong \Aut(T)^d\rtimes S_d=\Aut(T)\wr S_d.
\]

By definition,
\[
    \Hol(G)=G\rtimes \Aut(G)
    =
    T^d\rtimes(\Aut(T)\wr S_d)
    =
    (T\rtimes \Aut(T))\wr S_d.
\]
For nonabelian simple \(T\),
\[
    \Aut(T)=\Inn(T)\rtimes \Out(T),
    \qquad
    \Inn(T)\cong T.
\]
In the holomorph action, the right regular copy \(\hat G\) and the inner
automorphism copy \(\check G\) centralize each other, so
\[
    \Hol(G)
    =
    (\check G\times \hat G).(\Out(T)\wr S_d).
\]
When \(d=1\), this gives
\[
    \Hol(T)=T\rtimes\Aut(T)
    =
    (\check T\times \hat T).\Out(T).
\]
\end{proof}

\begin{exercise}
Let \(G\) be a finite group, and \(M \lhd G\). Prove that \(G\) acts transitively
on \(\operatorname{Syl}_p(M)\) where \(p\) is a prime divisor of \(|M|\), and
\[
    G = M N_G(M \cap P)
\]
for \(P \in \operatorname{Syl}_p(G)\).
\end{exercise}

\begin{proof}
Since \(M\trianglelefteq G\), conjugation by \(G\) sends Sylow \(p\)-subgroups
of \(M\) to Sylow \(p\)-subgroups of \(M\). Thus \(G\) acts on
\(\Syl_p(M)\). The subgroup \(M\) acts transitively on \(\Syl_p(M)\) by the
Sylow conjugacy theorem, hence the larger group \(G\) also acts transitively.

It remains to prove the factorization. Since \(P\in \Syl_p(G)\) and
\(M\trianglelefteq G\), the subgroup
\[
    M\cap P
\]
is a Sylow \(p\)-subgroup of \(M\). Indeed, the image \(MP/M\cong P/(M\cap P)\)
is a \(p\)-subgroup of \(G/M\), so \(M\cap P\) contains the full \(p\)-part of
\(M\).

Apply the Frattini argument to \(M\trianglelefteq G\) and the Sylow subgroup
\(M\cap P\in \Syl_p(M)\). We obtain
\[
    G=M N_G(M\cap P).
\]
\end{proof}

\begin{exercise}
Let \(H,K\) be subgroups of a group \(G\). Prove that \(G = HK\) if and only if
\[
    |H||K| = |G||H \cap K|.
\]
\end{exercise}

\begin{proof}
For finite subgroups \(H,K\leq G\), the multiplication map
\[
    H\times K\longrightarrow HK,
    \qquad
    (h,k)\mapsto hk
\]
has fibers of size \(|H\cap K|\). Hence
\[
    |HK|=\frac{|H||K|}{|H\cap K|}.
\]
Therefore \(G=HK\) if and only if \(|HK|=|G|\), which is equivalent to
\[
    |H||K|=|G||H\cap K|.
\]
\end{proof}

\begin{exercise}
If \(H,K < G\) such that \(|G:H|\) and \(|G:K|\) are coprime, prove
\(G = HK\).
\end{exercise}

\begin{proof}
Put
\[
    m=[G:H],
    \qquad
    n=[G:K],
    \qquad
    \gcd(m,n)=1.
\]
Since
\[
    [G:H\cap K]=[G:H][H:H\cap K],
\]
the integer \([G:H]\) divides \([G:H\cap K]\). Similarly, \([G:K]\) divides
\([G:H\cap K]\). Hence \(mn\mid [G:H\cap K]\). On the other hand,
\[
    [G:H\cap K]\leq [G:H][G:K]=mn.
\]
Therefore
\[
    [G:H\cap K]=mn.
\]
Thus
\[
    |H\cap K|=\frac{|G|}{mn}.
\]
Consequently
\[
    |H||K|
    =
    \frac{|G|}{m}\frac{|G|}{n}
    =
    |G|\frac{|G|}{mn}
    =
    |G||H\cap K|.
\]
By the previous exercise, \(G=HK\).
\end{proof}

\begin{exercise}
Let \(G = A_n \times A_n\), with \(n \geq 5\) odd, be a primitive permutation
group of HS type. Prove that \(G\) has a regular subgroup which is isomorphic to
\(\mathbb{Z}_n \times A_{n-1}\).
\end{exercise}

\begin{proof}
Let \(T=A_n\). In the HS action, \(T\times T\) acts on the set \(T\) by left and
right multiplication. The stabilizer of the identity element is the diagonal
subgroup
\[
    \{(t,t)\mid t\in T\}.
\]
Let
\[
    c=(1\,2\,\cdots\, n).
\]
Since \(n\) is odd, \(c\in A_n\), and \(\langle c\rangle\cong C_n\). Let
\[
    H=A_{n-1}
\]
be the stabilizer of the point \(n\) in the natural action of \(A_n\). Then
\[
    |\langle c\rangle|\,|H|
    =
    n\cdot \frac{(n-1)!}{2}
    =
    \frac{n!}{2}
    =
    |A_n|.
\]
Moreover
\[
    \langle c\rangle\cap H=1,
\]
because a nontrivial power of the \(n\)-cycle \(c\) fixes no point.

Now consider
\[
    R=\langle c\rangle\times H\leq A_n\times A_n.
\]
Its order is \(|T|\). Its stabilizer of \(1\in T\) is
\[
    R\cap \{(t,t)\mid t\in T\}
    =
    \{(1,1)\},
\]
because \(\langle c\rangle\cap H=1\). Hence \(R\) is semiregular of order
\(|T|\), and therefore regular. Thus \(G\) contains a regular subgroup
\[
    R\cong C_n\times A_{n-1}.
\]
\end{proof}

\begin{exercise}
State and prove Sylow's theorem.
\end{exercise}

\begin{proof}
Let \(G\) be finite and write
\[
    |G|=p^e m,
    \qquad
    \gcd(p,m)=1.
\]
Sylow's theorem consists of the following three statements.

\begin{enumerate}[label=\textup{(\arabic*)}]
    \item \(G\) has a subgroup of order \(p^e\).
    \item Any two subgroups of order \(p^e\) are conjugate. More generally, every
    \(p\)-subgroup of \(G\) is contained in a conjugate of a fixed Sylow
    \(p\)-subgroup.
    \item If \(n_p=|\Syl_p(G)|\), then
    \[
        n_p\mid m,
        \qquad
        n_p\equiv 1\pmod p.
    \]
\end{enumerate}

For existence, argue by induction on \(|G|\). Using the class equation, if
\(p\nmid |Z(G)|\), then one noncentral conjugacy class has size not divisible by
\(p\). Its centralizer is a proper subgroup whose order is still divisible by
\(p^e\), so induction gives a Sylow \(p\)-subgroup. If \(p\mid |Z(G)|\), then
Cauchy's theorem for the abelian group \(Z(G)\) gives a central subgroup
\(C\) of order \(p\). Induction applied to \(G/C\) lifts a subgroup of order
\(p^{e-1}\) to a subgroup of \(G\) of order \(p^e\).

For conjugacy, fix \(P\in \Syl_p(G)\) and let
\[
    \Delta=\{P^g\mid g\in G\}.
\]
Then
\[
    |\Delta|=[G:N_G(P)]\mid [G:P]=m,
\]
so \(p\nmid |\Delta|\). If \(H\leq G\) is a \(p\)-subgroup, let \(H\) act on
\(\Delta\) by conjugation. Since \(|\Delta|\) is not divisible by \(p\), some
\(P^g\) is fixed by \(H\). Hence
\[
    H\leq N_G(P^g).
\]
But \(P^g\trianglelefteq N_G(P^g)\) is the Sylow \(p\)-subgroup of its
normalizer, so \(H\leq P^g\). Taking \(H\) itself Sylow gives conjugacy of
Sylow \(p\)-subgroups.

Finally,
\[
    n_p=|\Syl_p(G)|=[G:N_G(P)]\mid [G:P]=m.
\]
Let \(P\) act on \(\Syl_p(G)\) by conjugation. The only fixed point is \(P\)
itself; all other orbits have size divisible by \(p\). Therefore
\[
    n_p\equiv 1\pmod p.
\]
\end{proof}

\begin{exercise}
Show that a group of order \(351 = 3^3 \cdot 13\) has a normal Sylow subgroup.
\end{exercise}

\begin{proof}
Let \(n_{13}\) be the number of Sylow \(13\)-subgroups. Then
\[
    n_{13}\mid 27,
    \qquad
    n_{13}\equiv 1\pmod {13}.
\]
Thus
\[
    n_{13}=1
    \quad\text{or}\quad
    n_{13}=27.
\]
If \(n_{13}=1\), the Sylow \(13\)-subgroup is normal.

Assume \(n_{13}=27\). Distinct Sylow \(13\)-subgroups intersect trivially, so
they contain
\[
    27(13-1)=324
\]
nonidentity elements. This leaves
\[
    351-324=27
\]
elements. These remaining elements must form a single Sylow \(3\)-subgroup.
Hence the Sylow \(3\)-subgroup is unique and normal. In either case \(G\) has a
normal Sylow subgroup.
\end{proof}

\begin{exercise}
Find two different Sylow \(2\)-subgroups of \(S_5\) and find an element that
conjugates one to the other.
\end{exercise}

\begin{proof}
A Sylow \(2\)-subgroup of \(S_5\) has order \(8\). Let
\[
    P=\langle (1\,2\,3\,4), (1\,3)\rangle.
\]
This is a dihedral group of order \(8\), acting on \(\{1,2,3,4\}\) and fixing
\(5\). Let
\[
    g=(4\,5).
\]
Then
\[
    P^g
    =
    \langle (1\,2\,3\,5),(1\,3)\rangle
    =
    Q.
\]
Thus \(P\) and \(Q\) are two different Sylow \(2\)-subgroups of \(S_5\), and
\((4\,5)\) conjugates \(P\) to \(Q\).
\end{proof}

\begin{exercise}
Prove that a group of order \(255 = 3 \cdot 5 \cdot 17\) is abelian.
\end{exercise}

\begin{proof}
Let \(P\in \Syl_{17}(G)\). Then
\[
    n_{17}\mid 15,
    \qquad
    n_{17}\equiv 1\pmod {17}.
\]
Hence \(n_{17}=1\), so \(P\trianglelefteq G\). Since \(P\cong C_{17}\),
conjugation gives
\[
    G/C_G(P)\lesssim \Aut(P)\cong C_{16}.
\]
But \(|G|=255\) is coprime to \(16\), so \(C_G(P)=G\). Thus
\[
    P\leq Z(G).
\]

Now
\[
    |G/P|=15.
\]
Every group of order \(15\) is cyclic, because the Sylow \(5\)-subgroup is
normal and the Sylow \(3\)-subgroup is also normal. Hence \(G/P\) is cyclic.
Since \(P\leq Z(G)\), the quotient \(G/Z(G)\) is a quotient of the cyclic group
\(G/P\), and is therefore cyclic. If \(G/Z(G)\) is cyclic, then \(G\) is
abelian. Hence \(G\) is abelian.
\end{proof}

\begin{exercise}
Let \(p,q,r\) be distinct primes with \(p > q > r\). Classify groups of order
\(pqr\).
\end{exercise}

\begin{proof}
The order \(pqr\) is square-free. By Holder's theorem for groups of square-free
order, every group of order \(pqr\) is metacyclic. More explicitly, every such
group has a presentation
\[
    G(d,e,k)
    =
    \langle x,y\mid x^d=y^e=1,\ y^{-1}xy=x^k\rangle,
\]
where
\[
    de=pqr,\qquad \gcd(d,e)=1,
\]
and
\[
    k\in(\mathbb Z/d\mathbb Z)^\times,
    \qquad
    k^e\equiv 1\pmod d.
\]
Thus
\[
    G\cong C_d\rtimes C_e,
\]
where the action of \(C_e\) on \(C_d\) is multiplication by \(k\).

Conversely, every such presentation defines a group of order \(pqr\). Two
presentations with the same \(d,e\) give the same group precisely up to replacing
the generator \(y\) of \(C_e\) by \(y^u\), where
\[
    u\in(\mathbb Z/e\mathbb Z)^\times;
\]
this replaces \(k\) by \(k^u\).

In concrete terms, the abelian group
\[
    C_p\times C_q\times C_r\cong C_{pqr}
\]
always occurs. Nonabelian groups occur exactly when one prime divisor of \(e\)
acts nontrivially on one prime divisor of \(d\). Such an action from a prime
\(s\) to a prime \(t\) exists exactly when
\[
    s\mid t-1,
\]
because
\[
    \Aut(C_t)\cong C_{t-1}.
\]
Thus the above semidirect products, with all allowed actions
\[
    r\mid q-1,\qquad r\mid p-1,\qquad q\mid p-1,
\]
and their compatible simultaneous actions, give all groups of order \(pqr\).
\end{proof}

\begin{exercise}
Assume that \(M \lhd G\). Show that
\[
    G/C_G(M) \lesssim \operatorname{Aut}(M).
\]
\end{exercise}

\begin{proof}
Let \(G\) act on \(M\) by conjugation. This gives a homomorphism
\[
    \theta:G\longrightarrow \Aut(M),
    \qquad
    \theta(g)(m)=g^{-1}mg.
\]
The kernel consists exactly of the elements of \(G\) that centralize every
element of \(M\):
\[
    \Ker(\theta)=C_G(M).
\]
By the first isomorphism theorem,
\[
    G/C_G(M)\cong \theta(G)\leq \Aut(M).
\]
\end{proof}

\begin{exercise}
Let \(G\) be a group of order \(72\), and assume that the center \(Z(G) = 1\).
Prove that \(G\) is isomorphic to one of the following groups:
\[
    \mathbb{Z}_3^2:\mathbb{Z}_8,\quad
    \mathbb{Z}_3^2:\mathrm{Q}_8,\quad
    \mathbb{Z}_3^2:\mathrm{D}_8,\quad
    (A_4 \times \mathbb{Z}_3):\mathbb{Z}_2.
\]
\end{exercise}

\begin{proof}
Since
\[
    |G|=72=2^3\cdot 3^2,
\]
the group \(G\) is solvable by Burnside's \(p^a q^b\)-theorem. Hence every
minimal normal subgroup of \(G\) is elementary abelian. The possible minimal
normal subgroups have order
\[
    2,\ 2^2,\ 2^3,\ 3,\ 3^2.
\]

The condition \(Z(G)=1\) eliminates the cases in which a minimal normal subgroup
is centralized by all of \(G\). The remaining calculation is the standard
linear action calculation on an elementary abelian minimal normal subgroup. In
the main case one obtains a self-centralizing normal subgroup
\[
    N\cong C_3^2.
\]
Then \(G/N\) has order \(8\), and the conjugation action embeds
\[
    G/N\hookrightarrow \Aut(N)\cong \GL(2,3).
\]
The subgroups of order \(8\) in \(\GL(2,3)\) which give centerless semidirect
products are, up to conjugacy,
\[
    C_8,\qquad Q_8,\qquad D_8.
\]
These give
\[
    C_3^2:C_8,\qquad
    C_3^2:Q_8,\qquad
    C_3^2:D_8.
\]

There is one further reducible action possibility. In that case the normal
\(3\)-subgroup splits into two \(1\)-dimensional constituents, and the induced
action interchanges the corresponding complements. This gives
\[
    (A_4\times C_3):C_2.
\]
These four possibilities are pairwise nonisomorphic and each has trivial
center. Therefore any group of order \(72\) with trivial center is isomorphic
to one of the four groups listed.
\end{proof}

\begin{exercise}
Let \(T\) be a simple group which is \(2\)-transitive on \(\Delta\), and let
\(G = T \wr S_3\) acts on \(\Omega = \Delta^3\) in product action. Prove that
\(G\) is of rank \(4\).
\end{exercise}

\begin{proof}
Fix \(\delta\in\Delta\), and put
\[
    \omega=(\delta,\delta,\delta)\in\Omega.
\]
The stabilizer \(G_\omega\) is
\[
    T_\delta^3\rtimes S_3.
\]
For
\[
    \alpha=(\alpha_1,\alpha_2,\alpha_3)\in\Delta^3,
\]
define the weight
\[
    \operatorname{wt}(\alpha)
    =
    |\{i\mid \alpha_i\neq \delta\}|.
\]
The point stabilizer \(G_\omega\) preserves this weight, so there are at least
four suborbits, corresponding to weights \(0,1,2,3\).

Since \(T\) is \(2\)-transitive, \(T_\delta\) is transitive on
\(\Delta\setminus\{\delta\}\). Hence two points of \(\Delta^3\) with the same
weight are conjugate under \(G_\omega\): the factors \(T_\delta^3\) change the
non-\(\delta\) entries, and \(S_3\) permutes their positions. Thus the four
weight classes are exactly the four \(G_\omega\)-orbits. Therefore
\[
    \rank(G)=4.
\]
\end{proof}

\begin{exercise}
Let \(T = A_n\), and let
\[
    G = T:S_n = (T \times T):S_2
\]
be a primitive permutation group of type HS. Find \(\operatorname{rank}(G)\).
\end{exercise}

\begin{proof}
In the HS action, identify the point set with \(T=A_n\). The point stabilizer is
\[
    \Aut(A_n)\cong S_n
\]
acting on \(A_n\) by conjugation. Hence the rank is the number of
\(S_n\)-conjugacy classes contained in \(A_n\).

The \(S_n\)-conjugacy classes are parametrized by partitions \(\lambda\vdash n\),
or equivalently by cycle type. A permutation of cycle type \(\lambda\) is even
if and only if
\[
    n-\ell(\lambda)
\]
is even, where \(\ell(\lambda)\) is the number of parts of \(\lambda\). Thus
\[
    \rank(G)
    =
    \#\{\lambda\vdash n\mid n-\ell(\lambda)\equiv 0\pmod 2\}.
\]
Equivalently, if \(n\) is odd, this is the number of partitions of \(n\) with
an odd number of parts.
\end{proof}

\begin{exercise}
Let \(T = A_5\), and \(N = T^k\). Determine the orbits of
\(\operatorname{Aut}(N)\) acting on the elements of \(N\).
\end{exercise}

\begin{proof}
Since
\[
    \Aut(A_5)\cong S_5,
\]
the \(\Aut(A_5)\)-orbits on \(A_5\) are determined by element order:
\[
    1,\qquad 2,\qquad 3,\qquad 5.
\]
Moreover
\[
    \Aut(N)=\Aut(A_5^k)\cong \Aut(A_5)\wr S_k.
\]
The base group \(\Aut(A_5)^k\) acts independently on the coordinates, and the
top group \(S_k\) permutes the coordinates. Therefore an element
\[
    (x_1,\ldots,x_k)\in A_5^k
\]
is determined up to \(\Aut(N)\)-orbit by the numbers of coordinates of each of
the four possible types:
\[
    |x_i|=1,\quad |x_i|=2,\quad |x_i|=3,\quad |x_i|=5.
\]
Thus the orbits are indexed by quadruples
\[
    (m_1,m_2,m_3,m_5)
\]
of nonnegative integers satisfying
\[
    m_1+m_2+m_3+m_5=k.
\]
In particular, the number of orbits is
\[
    \binom{k+3}{3}.
\]
\end{proof}

\begin{exercise}
Let \(G\) be a primitive permutation group of rank \(3\). Prove that \(G\) is
affine, almost simple, or \(T \wr S_2\) acting on \(\Delta^2\) in product action
with \(T\) being \(2\)-transitive on \(\Delta\).
\end{exercise}

\begin{proof}
A primitive permutation group is quasiprimitive. Apply the quasiprimitive
O'Nan--Scott analysis.

If \(G\) is of affine type, we are in the first case. If \(G\) is almost simple,
we are in the second case. It remains to exclude the other types or reduce them
to product action.

The notes prove that a group of holomorph simple type has rank at least \(4\),
and that compound regular nonabelian types have even larger rank. Diagonal type
also has rank at least \(4\). Hence none of these types can occur when
\(\rank(G)=3\).

Thus the only remaining possibility is product action:
\[
    G\leq T\wr S_k
    \quad\text{on}\quad
    \Omega=\Delta^k.
\]
The point stabilizer preserves the weight of a point, so
\[
    \rank(G)\geq k+1.
\]
Since \(\rank(G)=3\), we must have \(k=2\). Equality in the weight bound occurs
exactly when \(T\) is \(2\)-transitive on \(\Delta\). Therefore the product
action case is precisely
\[
    T\wr S_2
    \quad\text{on}\quad
    \Delta^2,
\]
with \(T\) \(2\)-transitive.
\end{proof}

\begin{exercise}
A transitive permutation group on \(\Omega\) is called bi-quasiprimitive if each
of its non-trivial normal subgroups has at most two orbits on \(\Omega\), and
there is one normal subgroup which has two orbits.
\begin{enumerate}[label=(\alph*)]
    \item Give examples of bi-quasiprimitive permutation groups.
    \item Prove that, a transitive permutation group is bi-quasiprimitive if and
    only if each of its minimal normal subgroups has exactly two orbits.
    \item How many minimal normal subgroups does a bi-quasiprimitive group have?
\end{enumerate}
\end{exercise}

\begin{proof}
\((a)\) Let \(T\) be a nonabelian simple transitive permutation group on
\(\Delta\). Put
\[
    G=(T\times T)\rtimes \langle \tau\rangle,
\]
where \(\tau\) interchanges the two factors, and let \(G\) act on the disjoint
union
\[
    \Omega=\Delta_1\sqcup \Delta_2
\]
by letting the two copies of \(T\) act on the two copies of \(\Delta\), while
\(\tau\) interchanges \(\Delta_1\) and \(\Delta_2\). Then
\[
    N=T\times T
\]
is normal in \(G\) and has exactly two orbits on \(\Omega\). The transitivity of
\(G\) comes from the element \(\tau\). This gives a basic family of
bi-quasiprimitive groups.

\((b)\) Suppose first that \(G\) is bi-quasiprimitive and let \(M\) be a minimal
normal subgroup of \(G\). By definition, \(M\) has at most two orbits. Since
there is a normal subgroup \(N\) with two orbits, a minimal normal subgroup
contained in \(N\) has two orbits. If another minimal normal subgroup were
transitive, then it would centralize this intransitive minimal normal subgroup,
forcing the two \(N\)-orbits to be blocks fixed by a transitive normal subgroup,
a contradiction. Hence every minimal normal subgroup has exactly two orbits.

Conversely, assume that every minimal normal subgroup of \(G\) has exactly two
orbits. Let \(1\neq N\lhd G\). Then \(N\) contains a minimal normal subgroup
\(M\). The \(N\)-orbits are unions of the two \(M\)-orbits, so \(N\) has at most
two orbits. Since the minimal normal subgroups themselves have two orbits, at
least one nontrivial normal subgroup has two orbits. Thus \(G\) is
bi-quasiprimitive.

\((c)\) A bi-quasiprimitive group has either one or two minimal normal
subgroups. Indeed, distinct minimal normal subgroups centralize one another,
and the product of any two is again normal with at most two orbits. Three
distinct minimal normal subgroups would force too many independent orbit
systems. Hence the number is \(1\) or \(2\).
\end{proof}

\begin{exercise}
For a transitive permutation group \(G\) on \(\Omega\) and \(\omega \in \Omega\),
an orbit of the stabilizer \(G_\omega\) on \(\Omega \setminus \{\omega\}\) is
called a suborbit of \(G\).
\begin{enumerate}[label=(\alph*)]
    \item Let \(G\) be a primitive permutation group which has a suborbit of
    length \(2\). Prove that \(G\) is a dihedral group of order \(2p\) with
    \(p\) prime.
    \item Let \(G \leq \operatorname{Sym}(\Omega)\) be primitive which has a
    suborbit of length \(3\) or \(4\). Prove that the stabilizer \(G_\omega\) is
    a \(\{2,3\}\)-group.
    \item Let \(G \leq \operatorname{Sym}(\Omega)\) be a \(p\)-group with
    \(p\) prime. Show that each suborbit of \(G\) has length being a
    \(p\)-power.
\end{enumerate}
\end{exercise}

\begin{proof}
\((a)\) Let \(\Delta\) be a suborbit of length \(2\), and form the corresponding
orbital graph with vertex set \(\Omega\). Since \(G\) is primitive, this graph
is connected. Its valency is \(2\), so it is a cycle. A primitive group acting
on a cycle has prime degree, say \(p\), and the full transitive group generated
by the rotations and a reflection is dihedral of order \(2p\). Hence
\[
    G\cong D_{2p}.
\]

\((b)\) This is the small-suborbit form of Wielandt's theorem. Let \(\Delta\) be
a suborbit of length \(3\) or \(4\). The action of \(G_\omega\) on the connected
orbital graph generated by \(\Delta\) is faithful once its action on the first
neighborhood is known. Hence \(G_\omega\) embeds into an iterated subgroup
generated from the local action on \(\Delta\). Since
\[
    |\Sym(3)|=6
    \qquad\text{and}\qquad
    |\Sym(4)|=24,
\]
only the primes \(2\) and \(3\) can divide \(|G_\omega|\). Thus \(G_\omega\) is
a \(\{2,3\}\)-group.

\((c)\) If \(G\) is a \(p\)-group, then every point stabilizer \(G_\omega\) is
also a \(p\)-group. A suborbit is an orbit of \(G_\omega\), so its length divides
\[
    |G_\omega|.
\]
Therefore each suborbit length is a power of \(p\).
\end{proof}

\begin{exercise}
A group is called generalised Hamiltonian if each subgroup is not core-free. For
example, a group \(A \times \mathrm{Q}_{2^n}\) where \(A\) is abelian and
\(\mathrm{Q}_{2^n}\) is generalised quaternion is a generalised Hamiltonian
group.
\begin{enumerate}[label=(\alph*)]
    \item Prove that an abelian transitive permutation group is regular.
    \item Show that a transitive permutation Hamiltonian group is regular.
    \item Let \(G\) be a permutation group which contains a transitive dihedral
    subgroup. Prove that \(G\) contains a cyclic regular subgroup or a dihedral
    regular subgroup.
\end{enumerate}
\end{exercise}

\begin{proof}
\((a)\) Let \(A\leq \Sym(\Omega)\) be abelian and transitive. For
\(\omega\in\Omega\), the stabilizer \(A_\omega\) is normal in \(A\), because
\(A\) is abelian. Hence \(A_\omega\) fixes every point:
\[
    A_\omega\leq \bigcap_{\alpha\in\Omega}A_\alpha.
\]
The action is faithful, so the kernel is \(1\). Thus \(A_\omega=1\), and \(A\)
is regular.

\((b)\) Let \(G\leq \Sym(\Omega)\) be transitive and Hamiltonian in the stated
sense. The core of a point stabilizer is the kernel of the action:
\[
    \Core_G(G_\omega)=\bigcap_{g\in G}G_\omega^g=1.
\]
If \(G_\omega\neq 1\), then \(G_\omega\) would be a core-free subgroup,
contrary to the Hamiltonian assumption. Hence \(G_\omega=1\), and \(G\) is
regular.

\((c)\) Let
\[
    D=\langle r,s\mid r^m=s^2=1,\ srs=r^{-1}\rangle
\]
be a transitive dihedral subgroup of \(G\). Its rotation subgroup
\[
    R=\langle r\rangle
\]
has at most two orbits. If \(R\) is transitive, then by part \((a)\) it is
regular, so \(G\) contains a cyclic regular subgroup. If \(R\) has two orbits,
then the faithful transitive action of \(D\) has trivial point stabilizer, and
so \(D\) itself is regular. Thus \(G\) contains a cyclic regular subgroup or a
dihedral regular subgroup.
\end{proof}

\begin{exercise}
Let \(G\) be a finite solvable group whose Sylow subgroups are cyclic or
dihedral. Describe the structure of \(G\).
\end{exercise}

\begin{proof}
Let
\[
    F=F(G)
\]
be the Fitting subgroup. Since \(G\) is solvable, the notes give
\[
    F=\prod_{p\mid |G|}O_p(G)
\]
and
\[
    C_G(F)\leq F.
\]
For odd \(p\), a Sylow \(p\)-subgroup cannot be dihedral, so it is cyclic.
Hence \(O_p(G)\) is cyclic for every odd \(p\). The \(2\)-part \(O_2(G)\) is a
subgroup of a cyclic or dihedral Sylow \(2\)-subgroup, and is therefore cyclic
or dihedral.

Thus \(F(G)\) is a direct product of cyclic odd-order normal Sylow cores and a
cyclic or dihedral normal \(2\)-core. Moreover the conjugation action gives
\[
    G/C_G(F)\lesssim \Aut(F),
\]
and since \(C_G(F)\leq F\), the quotient \(G/F\) embeds in a quotient of
\(\Aut(F)\). Therefore \(G\) is an extension of such a nilpotent group \(F\) by
a subgroup of \(\Aut(F)\). In the special case where all Sylow subgroups are
cyclic, \(F\) is cyclic and \(G/F\) is cyclic, so \(G\) is metacyclic.
\end{proof}

\begin{exercise}
Let \(G\) be a finite solvable group. If \(F(G)\) is a \(p\)-group, then
\(F(G/F(G))\) is a \(p'\)-group.
\end{exercise}

\begin{proof}
Put \(F=F(G)\), and suppose for contradiction that
\[
    F(G/F)
\]
has a nontrivial normal \(p\)-subgroup. Let \(Y/F\) be its full preimage in
\(G/F\). Then \(Y\lhd G\), and \(Y/F\) is a \(p\)-group. Since \(F\) is also a
\(p\)-group, \(Y\) is a normal \(p\)-subgroup of \(G\). Hence \(Y\) is nilpotent
and normal, so by the maximality of the Fitting subgroup,
\[
    Y\leq F.
\]
This contradicts \(Y/F\neq 1\). Therefore the Fitting subgroup of \(G/F\) has
no nontrivial \(p\)-part; equivalently,
\[
    F(G/F(G))
\]
is a \(p'\)-group.
\end{proof}

\begin{exercise}
Let \(G\) be solvable and \(\Phi(G) = 1\). Prove the following statements.
\begin{enumerate}[label=(\alph*)]
    \item \(F(G)\) is the only maximal abelian normal subgroup of \(G\).
    \item If \(G\) has a unique minimal normal subgroup \(N\), then
    \(F(G) = N\).
\end{enumerate}
\end{exercise}

\begin{proof}
For solvable groups with \(\Phi(G)=1\), Gasch\"utz's theorem gives that
\[
    F(G)
\]
is the direct product of the minimal normal subgroups of \(G\). In particular,
\(F(G)\) is abelian.

\((a)\) Let \(A\lhd G\) be abelian. Then \(A\) is nilpotent and normal, so
\[
    A\leq F(G).
\]
Since \(F(G)\) itself is abelian and normal, it contains every abelian normal
subgroup of \(G\). Therefore \(F(G)\) is the unique maximal abelian normal
subgroup of \(G\).

\((b)\) If \(G\) has a unique minimal normal subgroup \(N\), then the direct
product of all minimal normal subgroups is just \(N\). Hence
\[
    F(G)=N.
\]
\end{proof}

\begin{exercise}
Let \(G\) be a group of order \(p^3\). Determine \(G'\), \(Z(G)\), and
\(\Phi(G)\).
\end{exercise}

\begin{proof}
If \(G\) is abelian, then \(G'=1\) and \(Z(G)=G\). The three abelian cases are:
\[
    G\cong C_{p^3},
    \qquad
    G\cong C_{p^2}\times C_p,
    \qquad
    G\cong C_p^3.
\]
For a finite abelian \(p\)-group,
\[
    \Phi(G)=pG.
\]
Thus
\[
    \Phi(C_{p^3})\cong C_{p^2},
\]
\[
    \Phi(C_{p^2}\times C_p)\cong C_p,
\]
and
\[
    \Phi(C_p^3)=1.
\]

If \(G\) is nonabelian of order \(p^3\), then the usual \(p^3\)-classification
gives
\[
    |Z(G)|=p,
    \qquad
    G'=Z(G)\cong C_p.
\]
Also, for a finite \(p\)-group,
\[
    \Phi(G)=G'G^p.
\]
In each nonabelian group of order \(p^3\), this equals the central subgroup of
order \(p\). Hence
\[
    G'=Z(G)=\Phi(G)\cong C_p
\]
in the nonabelian case.
\end{proof}

\begin{exercise}
Let \(A = \langle a \rangle = \mathbb{Z}_4\) and \(B = \mathrm{Q}_8\) with
\(Z(B) = \langle c \rangle\). Let \(G = A \circ B\) be the central product of
\(A\) and \(B\), namely,
\[
    G = (A \times B)/\langle a^2 c \rangle.
\]
List all the involutions of \(G\), namely, elements of \(G\) of order \(2\).
\end{exercise}

\begin{proof}
Write
\[
    B=Q_8=\{1,c,\pm i,\pm j,\pm k\},
    \qquad c=-1.
\]
In the quotient
\[
    G=(A\times B)/\langle a^2c\rangle,
\]
we identify
\[
    a^2=c.
\]
Every element of \(G\) has a representative
\[
    a^\epsilon b,
    \qquad
    \epsilon\in\{0,1\},
    \qquad
    b\in B.
\]
Since \(A\) and \(B\) commute in the central product,
\[
    (a^\epsilon b)^2=a^{2\epsilon}b^2.
\]
If \(\epsilon=0\), the only nonidentity element with square \(1\) is
\[
    c=a^2.
\]
If \(\epsilon=1\), then \((ab)^2=1\) exactly when \(b^2=c\), that is, when
\[
    b\in \{\pm i,\pm j,\pm k\}.
\]
Therefore the involutions of \(G\) are
\[
    a^2,
    \qquad
    ai,\ a(-i),\ aj,\ a(-j),\ ak,\ a(-k).
\]
\end{proof}

\begin{exercise}
State and prove Hall's theorems for solvable groups.
\end{exercise}

\begin{proof}
Let \(G\) be a finite solvable group and let
\[
    \pi\subseteq \pi(G).
\]
Hall's theorems for solvable groups say:
\begin{enumerate}[label=\textup{(\arabic*)}]
    \item \(G\) has a Hall \(\pi\)-subgroup.
    \item Any two Hall \(\pi\)-subgroups of \(G\) are conjugate.
    \item Every \(\pi\)-subgroup of \(G\) is contained in a Hall
    \(\pi\)-subgroup.
\end{enumerate}

For existence, use induction on \(|G|\). Let \(M\) be a minimal normal subgroup.
Since \(G\) is solvable,
\[
    M\cong C_p^d
\]
for some prime \(p\). If \(p\in\pi\), take the full preimage of a Hall
\(\pi\)-subgroup of \(G/M\). If \(p\notin\pi\), take the full preimage \(X\) of
a Hall \(\pi\)-subgroup of \(G/M\). If \(X<G\), apply induction inside \(X\).
If \(X=G\), then \(M\) has a complement in \(G\), and this complement is a Hall
\(\pi\)-subgroup.

For conjugacy, again use induction. Let \(H,K\) be Hall \(\pi\)-subgroups and
let \(M\lhd G\) be maximal normal with \(G/M\) cyclic of prime order. If that
prime is not in \(\pi\), then \(H,K\leq M\), so induction in \(M\) conjugates
them. If the prime is in \(\pi\), then \(HM=KM=G\), and induction first
conjugates \(H\cap M\) to \(K\cap M\). After this reduction, both \(H\) and
\(K\) lie in the normalizer of their common intersection, and induction or
Sylow conjugacy in the quotient gives \(H^g=K\).

Finally, if \(U\) is a \(\pi\)-subgroup, apply existence to a suitable normalizer
step and then conjugacy to place \(U\) inside one of the Hall \(\pi\)-subgroups.
This proves the containment statement.
\end{proof}

\begin{exercise}
Let \(G\) be a finite solvable group with
\[
    \pi(G) = \{p_1,\ldots,p_t\}.
\]
Let \(G_{p_i'}\) be a Hall \(p_i'\)-subgroup, with \(1 \leq i \leq t\), and let
\[
    \mathcal{K} = \{G_{p_1'},\ldots,G_{p_t'}\}.
\]
Prove
\begin{enumerate}[label=(\alph*)]
    \item for any \(I \subset \{1,\ldots,t\}\), the intersection
    \(\bigcap_{i\in I} G_{p_i'}\) is a Hall subgroup;
    \item \(G\) has Sylow \(p_i\)-subgroups \(G_{p_i}\) such that
    \(G_{p_i}G_{p_j} = G_{p_j}G_{p_i}\).
\end{enumerate}
\end{exercise}

\begin{proof}
\((a)\) For each \(i\), the index
\[
    [G:G_{p_i'}]
\]
is a power of \(p_i\). Therefore, for distinct \(i\), these indices are
pairwise coprime. By the coprime-index argument proved above, for any
\[
    I\subseteq \{1,\ldots,t\},
\]
we have
\[
    \left[G:\bigcap_{i\in I}G_{p_i'}\right]
    =
    \prod_{i\in I}[G:G_{p_i'}].
\]
Thus
\[
    \bigcap_{i\in I}G_{p_i'}
\]
has precisely the full Sylow parts for the primes not in \(I\), and none of the
Sylow parts for the primes in \(I\). Hence it is a Hall subgroup.

\((b)\) Define
\[
    G_{p_i}
    =
    \bigcap_{j\neq i}G_{p_j'}.
\]
By part \((a)\), this is a Hall \(\{p_i\}\)-subgroup, hence a Sylow
\(p_i\)-subgroup.

For \(i\neq j\), put
\[
    H_{ij}
    =
    \bigcap_{\ell\neq i,j}G_{p_\ell'}.
\]
Again by part \((a)\), \(H_{ij}\) is a Hall \(\{p_i,p_j\}\)-subgroup, and it
contains both \(G_{p_i}\) and \(G_{p_j}\). Since
\[
    |G_{p_i}G_{p_j}|=|G_{p_i}|\,|G_{p_j}|=|H_{ij}|,
\]
we get
\[
    G_{p_i}G_{p_j}=H_{ij}=G_{p_j}G_{p_i}.
\]
Thus the Sylow subgroups \(G_{p_i}\) can be chosen pairwise permutable.
\end{proof}
